\subsection{n-th SYMMETRIC POWER OF A MODULE AND SYMMETRIC MULTI-LINEAR MAPPINGS}

Let X, Y be two sets and n an integer \(\geqslant1\) . A symmetric mapping of \(X^{n}\) into Y is any mapping \(f: X^{n} \to Y\) such that, for every permutation \(\sigma \in S_{n}\) and every element \((x_{i}) \in X^{n}\) ,

\[
f (x _ {\sigma (1)}, x _ {\sigma (2)}, \dots , x _ {\sigma (n)}) = f (x _ {1}, x _ {2}, \dots , x _ {n}).\tag{5}
\]

As the transpositions which exchange two consecutive integers generate the group \(\mathfrak{S}_n\) (I, § 5, no. 7), it suffices that condition (5) hold when \(\sigma\) is such a transposition.

When Y is a module over a commutative ring A, clearly the set of symmetric mappings of \(X^{n}\) into Y is a submodule of the A-module \(Y^{X^{n}}\) of all mappings of \(X^{n}\) into Y.

PROPOSITION 6. Let A be a commutative ring and M and N two A-modules. If with every A-linear mapping \(g: \mathbf{S}^n(\mathbf{M}) \to \mathbf{N} (n \geqslant 1)\) is associated the n-linear mapping

\[
\left(x _ {1}, x _ {2}, \dots , x _ {n}\right) \mapsto g \left(x _ {1} x _ {2} \dots x _ {n}\right)\tag{6}
\]

(where on the right-hand side the product is taken in the algebra \(\mathbf{S}(\mathbf{M})\) ), a bijective A-linear mapping is obtained of the A-module \(\operatorname{Hom}_{\mathbf{A}}(\mathbf{S}^{n}(\mathbf{M}), \mathbf{N})\) onto the A-module of symmetric \(n\) -linear mappings of \(\mathbf{M}^{n}\) into \(\mathbf{N}\) .

Recall (II, §3, no. 9) that there is a canonical bijection of the A-module \(\operatorname{Hom}_{\mathbf{A}}(\mathbb{T}^{n}(\mathbf{M}), \mathbf{N})\) onto the A-module \(\mathcal{L}_{n}(\mathbf{M}, \ldots, \mathbf{M}; \mathbf{N})\) of all \(n\) -linear mappings of \(\mathbf{M}^{n}\) into \(\mathbf{N}\) obtained by associating with every A-linear mapping \(f: \mathbb{T}^{n}(\mathbf{M}) \to \mathbf{N}\) the \(n\) -linear mapping

\[
\bar {f} \colon (x _ {1}, x _ {2}, \dots , x _ {n}) \mapsto f (x _ {1} \otimes x _ {2} \otimes \dots \otimes x _ {n}).\tag{7}
\]

On the other hand, the A-linear mappings \(g \colon \mathbb{S}^n(\mathbf{M}) \to \mathbf{N}\) are in one-to-one correspondence with the A-linear mappings \(f \colon \mathbb{T}^n(\mathbf{M}) \to \mathbf{N}\) such that \(f\) is zero on \(\mathfrak{J}_n'\) , by associating with \(g\) the mapping \(f = g \circ p_n\) , where

\[
p _ {n} \colon \mathbf {T} ^ {n} (\mathbf {M}) \rightarrow \mathbf {S} ^ {n} (\mathbf {M}) = \mathbf {T} ^ {n} (\mathbf {M}) / \mathfrak {J} _ {n} ^ {\prime}
\]

is the canonical homomorphism (II, § 2, no. 1, Theorem 1). But as \(\mathfrak{J}_n'\) is a linear combination of elements of the form

\[
\left(u _ {1} \otimes u _ {2} \otimes \dots \otimes u _ {p}\right) \otimes (x \otimes y - y \otimes x) \otimes \left(v _ {1} \otimes \dots \otimes v _ {n - p - 2}\right)
\]

\((x, y, u_{i}, v_{j} \text{ in } \mathbf{M})\) , to say that the function f is of the form \(g \circ p_{n}\) means that the corresponding n-linear function \(\bar{f}\) satisfies the relation

\[
\bar {f} (u _ {1}, \dots , u _ {p}, x, y, v _ {1}, \dots , v _ {n - p - 2}) = \bar {f} (u _ {1}, \dots , u _ {p}, y, x, v _ {1}, \dots , v _ {n - p - 2});
\]

in other words, by what has been seen above, this means that \(\bar{f}\) is symmetric; whence the proposition, taking account of the fact that

\[
p _ {n} (x _ {1} \otimes x _ {2} \otimes \dots \otimes x _ {n}) = x _ {1} x _ {2} \dots x _ {n}
\]

with the \(x_{i} \in M\) .

The A-module \(\mathbf{S}^{n}(\mathbf{M})\) is called the n-th symmetric power of M. For every A-module homomorphism \(u: M \to N\) , the mapping \(\mathbf{S}^{n}(u): \mathbf{S}^{n}(\mathbf{M}) \to \mathbf{S}^{n}(\mathbf{N})\) which coincides with \(\mathbf{S}(u)\) on \(\mathbf{S}^{n}(\mathbf{M})\) is called the n-th symmetric power of u.

Remark. Let \(\sigma\) be a permutation in \(\mathfrak{S}_n\) ; as the mapping

\[
\left(x _ {1}, x _ {2}, \dots , x _ {n}\right) \mapsto x _ {\sigma^ {- 1} (1)} \otimes x _ {\sigma^ {- 1} (2)} \otimes \dots \otimes x _ {\sigma^ {- 1} (n)}
\]

of \(\mathbf{M}^n\) into \(\mathbf{T}^n (\mathbf{M})\) is A-multilinear, it may be written uniquely as

\[
(x _ {1}, \dots , x _ {n}) \mapsto u _ {\sigma} (x _ {1} \otimes x _ {2} \otimes \dots \otimes x _ {n}),
\]

where \(u_{\sigma}\) is an endomorphism of the A-module \(\mathsf{T}^n(\mathbf{M})\) , also denoted by \(z \mapsto \sigma. z\) . Clearly, if \(\sigma\) is the identity element of \(\mathfrak{S}_n\) , \(u_{\sigma}\) is the identity; on the other hand, writing \(y_i = x_{\sigma^{-1}(i)}\) , we obtain, for every permutation \(\tau \in \mathfrak{S}_n, y_{\tau^{-1}(i)} = x_{\sigma^{-1}(\tau^{-1}(i))}\) and hence \(\tau. (\sigma. z) = (\tau \sigma). z\) ; in other words, the A-module \(\mathsf{T}^n(\mathbf{M})\) is a left \(\mathfrak{S}_n\text{-set}\) under the operation \((\sigma, z) \mapsto \sigma. z\) (I, § 5, no. 1). The elements of \(\mathsf{T}^n(\mathbf{M})\) such that \(\sigma. z = z\) for all \(\sigma \in \mathfrak{S}_n\) are called (contravariant) symmetric tensors of order \(n\) ; they form a sub-A-module \(\mathsf{S}_n'(\mathbf{M})\) of \(\mathsf{T}^n(\mathbf{M})\) .

For all \(z \in \mathbf{T}^n(\mathbf{M})\) , we write \(s.z = \sum_{\sigma \in \mathfrak{S}_n} \sigma.z\) and call \(s.z\) the symmetrization of the tensor \(z\) ; clearly \(s.z\) is a symmetric tensor and therefore \(z \mapsto s.z\) is an endomorphism of \(\mathbf{T}^n(\mathbf{M})\) whose image \(\mathbf{S}_n''(\mathbf{M})\) is contained in \(\mathbf{S}_n'(\mathbf{M})\) ; in general, \(\mathbf{S}_n''(\mathbf{M}) \neq \mathbf{S}_n'( \mathbf{M})\) (Exercise 5). If \(z\) is a symmetric tensor, then \(s.z = n!z\) ; it follows that when \(n!\) is invertible in \(\mathbf{A}\) , the endomorphism \(z \mapsto (n!)^{-1}s.z\) is a projector of \(\mathbf{T}^n(\mathbf{M})\) (II, § 1, no. 8), whose image is \(\mathbf{S}_n'(\mathbf{M}) = \mathbf{S}_n''(\mathbf{M})\) ; moreover the kernel on this projector is just \(\mathfrak{J}_n'\) . For, obviously \(\sigma(\mathfrak{J}_n') \subset \mathfrak{J}_n'\) for all \(\sigma \in \mathfrak{S}_n\) and \(\mathfrak{J}_n'\) is by definition generated by the tensors \(z - \rho.z\) , where \(\rho\) is a transposition exchanging two consecutive numbers in {[}1, n{]}; also, if \(\sigma, \tau\) are two permutations in \(\mathfrak{S}_n\) , then \(z - (\sigma\tau).z = z - \sigma.z + \sigma.(z - \tau.z)\) , whence it follows (since every permutation in \(\mathfrak{S}_n\) is a product of transpositions exchanging two consecutive numbers) that \(z - \sigma.z \in \mathfrak{J}_n'\) for all \(z \in \mathbf{T}^n(\mathbf{M})\) and \(\sigma \in \mathfrak{S}_n\) . Therefore (always supposing that \(n!\) is invertible in \(\mathbf{A}\) ), it is seen that \(z - (n!)^{-1}s.z = \sum_{\sigma \in \mathfrak{S}_n}(n!)^{-1}(z - \sigma.z) \in \mathfrak{J}_n'\) for all \(z \in \mathbf{T}^n(\mathbf{M})\) , which proves our assertion.

When \(n!\) is invertible in A, the submodules \(S_n'(M)\) and \(\mathfrak{S}_n'\) of \(T^n(M)\) are therefore supplementary and the restriction to \(S_n'(M)\) of the canonical homomorphism \(T^n(M) \to S^n(M) = T^n(M)/\mathfrak{J}_n'\) is an A-module isomorphism, which allows us in the case envisaged to identify symmetric tensors of order \(n\) with the elements of the \(n\) -th symmetric power of M. Note however that this identification is not compatible with multiplication, the product (in \(T(M)\) ) of two symmetric tensors not being symmetric in general and not therefore having as image in \(S(M)\) the product of the images of the symmetric tensors considered.

